\section{Positive-Cone Gauges and Common Mode}
\label{app:cone}

Every signed vector $x$ admits nonnegative representatives
$u=(x^+,x^-)^{\trans}$ with $\Pi u=x$.  Representatives are not unique:
$(x^++a,x^-+a)$ has the same projection for every $a\ge0$.  This freedom is
the common-mode gauge.

Canonicalization in \eqref{eq:rail-canonicalization} removes exactly
$2\sum_i\min(p_i,n_i)$ units of two-rail mass.  The identity $K^2=K$ makes the
remaining state a fixed point of further annihilation, while $\Pi K=\Pi$ makes
the removed mass an exact common-mode ledger.  No signed Fourier information is
carried by that ledger.

For a stage $M$, $C(M)u$ is nonnegative and projects to $Mx$.  For two stages,
\begin{equation}
 \Pi C(B)C(A)u=BA\Pi u,
 \label{eq:cone-two-stage}
\end{equation}
even when $C(B)C(A)\ne C(BA)$.  Their difference maps every representative into
$\ker\Pi$.  Hence annihilating equal rails between stages selects a smaller
representative without changing any downstream signed observable.

For the backward gauge, positivity of $g_s$ follows whenever every input
coordinate influences at least one terminal coordinate.  Normalized Bruun
cells have this property.  Column stochasticity gives
\begin{equation}
 \bm{1}^{\trans}C(A_s)u=\bm{1}^{\trans}u,
\end{equation}
so transported mass is conserved exactly.  For the forward gauge, row
stochasticity makes constant potential an exact fixed point.  The two gauges
are dual boundary normalizations of the same positive cone operator.
